Table~\ref{tab:detection} gives the mean AUC of every layer for the face map label. On 10\,m block means most layers rank face cells well (interpolation {{VAL:auc_I10_BICUBIC}}, boosting {{VAL:auc_I10_GBM}}, the detector {{VAL:auc_I10_DET}}, slope {{VAL:auc_I10_SLOPE}}), and the real USGS 1/3 arc second product behaves almost like the 10\,m block mean (interpolation {{VAL:auc_R10_BICUBIC}}). At 30\,m every elevation layer loses skill. With the settings of our earlier benchmark, interpolation is best ({{VAL:auc_I30_BICUBIC}}), boosting ({{VAL:auc_I30_GBM}}) and the network ({{VAL:auc_I30_CNN}}) follow, and the detector ({{VAL:auc_I30_DET}}) beats only slope ({{VAL:auc_I30_SLOPE}}) and the aspect rule ({{VAL:auc_I30_ASPECT}}). The learned layers were held back by those settings: learned super resolution of the 30\,m surface reaches {{VAL:auc_I30_SR}} and the network trained four times longer {{VAL:auc_I30_CNNlong}}, the best elevation layers at 30\,m, and boosting with five more features, among them the interpolated slope, {{VAL:auc_I30_GBMplus}}, which supports the idea that interpolation helps through local curvature. Tuning boosting by nested leave one site out changes nothing ({{VAL:auc_I30_GBMtuned}}). The information is not gone at 10\,m: the largest of the nine 10\,m interpolation scores in a 30\,m cell reaches {{VAL:auc_I10to30_BICUBIC}}.

Imagery adds what elevation lacks. Boosting on terrain plus Sentinel-2 features reaches {{VAL:auc_I30_DEMplusS2}} at 30\,m ({{VAL:gain_s2_I30}} over terrain alone), and with NAIP {{VAL:auc_I30_DEMplusNAIP}}; imagery alone is weaker ({{VAL:auc_I30_IMGS2}} and {{VAL:auc_I30_IMGNAIP}}), and at 10\,m it adds nothing ({{VAL:auc_I10_DEMplusS2}}). Super resolution and boosting with Sentinel-2 are above interpolation at all six sites and the detector below it at all six, but with six sites none of these comparisons is significant after Holm correction (sign flip $p={{VAL:sf_p_SR}}$, adjusted {{VAL:sf_holm_SR}}).

The real 30\,m products are worse than the 30\,m block mean: on their own grids the best layer reaches {{VAL:auc_G30n_GBM}} on GLO-30 and {{VAL:auc_N30_GBM}} on NASADEM, which predates the lidar by about two decades. Interpolating GLO-30 to the lidar grid, as in our earlier runs, changes little (best {{VAL:auc_G30_CNN}}).

The label definition matters less than expected. With the earlier cell label (at least 10 steep posts) the 30\,m interpolation scores {{VAL:lab_LCELLn10_BICUBIC}} against {{VAL:auc_I30_BICUBIC}}; requiring 100 or 300 steep posts raises it to {{VAL:lab_LCELLn100_BICUBIC}} and {{VAL:lab_LCELLn300_BICUBIC}}, because large faces are easier to see; a 30\,degree threshold changes little ({{VAL:lab_LANY30deg_BICUBIC}}). The labels are spatially correlated over {{VAL:corr_range_min}} to {{VAL:corr_range_max}}\,m (the largest at Dumont, whose window holds rocky hills); with blocks of {{VAL:corr_block_m}}\,m instead of 300\,m the per site intervals widen, but slope and the detector stay significantly below interpolation at all six sites.
